\subsection{Predictions for Ultra-Diffuse Galaxies}
  \label{subsec:udg}

  Ultra-diffuse galaxies constitute extreme low-surface-density systems and
  probe the saturated regime across most of their spatial
  extent~\cite{vanDokkum2016UDG}.

  Quantitatively, in the saturated regime ($g_{\mathrm{N}} \ll a_\star$), the
  effective framework predicts that the asymptotic circular velocity is bounded
  by the finite response of the effective gravitational field.
  To leading order, this implies
  \begin{equation}
    v_{\infty}
    \;\sim\;
    \sqrt{R\, a_\star}\;
    \left[
      1
      -
      \mathcal{O}\!\left(\delta_\rho\right)
    \right] ,
    \label{eq:udg-velocity}
  \end{equation}
  where $R$ is the characteristic baryonic scale length of the system and
  $\delta_\rho = 1 - \rho_{\mathrm{void}}/\rho_{\mathrm{global}}$ is the local
  underdensity contrast.
  For ultra-diffuse galaxies in void environments
  ($\delta_\rho \approx 0.3$--$0.6$ and $R \sim 3$--$5$~kpc),
  this yields typical asymptotic velocities
  $v_{\infty} \lesssim 10$--$15~\mathrm{km\,s^{-1}}$.

  This prediction differs qualitatively from MOND, which enforces a universal
  relation $v_{\infty} \propto (M_{\mathrm{baryon}} a_0)^{1/4}$ independent of
  environment.
  In the present framework, the asymptotic velocity depends explicitly on the
  large-scale environment through $\delta_\rho$, leading to systematically lower
  velocities for ultra-diffuse galaxies in voids than for systems of comparable
  baryonic mass in denser regions.

  Observations reveal a wide diversity of inferred dynamical mass content among
  ultra-diffuse galaxies, ranging from apparently dark-matter-rich systems to
  galaxies consistent with little or no dark
  matter~\cite{vanDokkum2019DF2,ManceraPina2019UDG}.
  Within the present framework, this diversity arises naturally from differences
  in how deeply individual systems probe the saturated regime.

  In ultra-diffuse galaxies, the apparent mass discrepancy inferred from
  kinematics should therefore not be interpreted as evidence for an extended
  distribution of unseen matter.
  It reflects a geometric threshold effect.
  The baryonic configuration probes a regime in which further density dilution no
  longer increases the effective dynamical response.

  As a result, the inferred dynamical mass tracks the onset of saturation rather
  than the presence of an additional gravitating component.
  The apparent ``missing mass'' is thus an emergent consequence of limited
  descriptive resolution in low-density systems.

  This interpretation leads to a falsifiable prediction.
  If the inferred mass discrepancy originates from geometric saturation rather
  than from a material mass component, independent mass tracers should not reveal
  correspondingly extended matter distributions.

  In particular, weak lensing signals, satellite dynamics, and stellar velocity
  dispersion profiles are expected to systematically underpredict the dynamical
  mass inferred from kinematic measurements in the deepest saturation regime.
  A positive detection of massive, spatially extended halos in ultra-diffuse
  galaxies would directly falsify the present framework.

  Because ultra-diffuse galaxies probe the deepest saturation regime, they
  provide a particularly sensitive test of the universality of the saturation
  scale $a_\star$.
